\SourceLeaf{106}
\begin{center}\textbf{[THREE FRAGMENTS ON ELLIPTIC MODULAR FUNCTIONS.]}\end{center}
\begin{center}\rule{1.5cm}{.4pt}\end{center}
\begin{center}[1.]\end{center}

If
\[
\frac{n}{m}=\frac{\mu(1-2e^{-M\pi}+2e^{-4M\pi}-2e^{-9M\pi}+\cdots)^2}{\mu(1+2e^{-M\pi}+2e^{-4M\pi}+2e^{-9M\pi}+\cdots)^2},
\]
then in place of $M$ one may put
\[
\frac{1}{2ai}\StepDown{.75}{+\frac{1}{2bi}}\StepDown{1.5}{+\frac{1}{2ci}}\StepDown{2.25}{+\frac{1}{2di}}\StepDown{3}{+\frac{1}{2ei}}\StepDown{3.75}{+\text{etc.}+\frac{1}{M},}
\]
where $a,b,c,d,e$ etc. denote any odd number of real integers. Alternatively one may put
\[
\frac{pM+2qi}{r+2sMi}=\frac{M+2(qr-psMM)i}{rr+4ssMM},
\]
where $p,q,r,s$ may be any real integers for which the condition
\[
pr+4qs=1
\]
holds.

\begin{center}\rule{3cm}{.4pt}\end{center}
\SourceLeaf{107}
\begin{center}[2.]\end{center}
\begin{center}\textbf{THE REDUCTION OF $pM$, $qM$, $rM$ TO THE SIMPLEST FORM.}\end{center}

Let $M=\dfrac{\alpha+\beta i}{\delta-\gamma i}$, where $\alpha$, $\beta$, $\gamma$, $\delta$ are real integers and the numerator and denominator have no common factor. Set
\[
\begin{gathered}
\alpha\alpha+\beta\beta=A,\qquad \alpha\gamma+\beta\delta=B,\qquad \gamma\gamma+\delta\delta=C,\\
\alpha\delta-\beta\gamma=\BareRoot{(AC-BB)}=D.
\end{gathered}
\]
Find the simplest form of determinant $-DD$ that is equivalent to the form $(A,B,C)$ and denote it by $(a,b,c)$.

For functions of $M$, reduction to functions of
\[
\frac{D+bi}{a}
\]
is then possible. The algorithm is as follows
\[
\begin{array}{ccc}
\dfrac{D+Bi}{A}=M,&DD+BB=AA',&B+B'=hA',\\[4pt]
\dfrac{D+B'i}{A'}=M',&DD+B'B'=A'A'',&B'+B''=h'A'',\\[4pt]
\dfrac{D+B''i}{A''}=M'',&DD+B''B''=A''A''',&B''+B'''=h''A''',
\end{array}
\]
\begin{center}.\quad .\quad .\quad .\quad .\quad .\quad .\quad .\quad .\quad .\quad .\quad .\quad .\quad .\quad .\quad .\quad .\quad .\quad .\quad .\quad .\end{center}
\[
\begin{array}{rl}
\BareRoot{\dfrac{D+Bi}{A}}\cdot pM=pM'&\text{for even }h,\\[4pt]
{}=qM'&\text{for odd }h,\\[5pt]\hline\\[-6pt]
\varepsilon^h\BareRoot{\dfrac{D+Bi}{A}}\cdot qM=rM',&\varepsilon=\sqrt{i},\\[5pt]\hline\\[-6pt]
\BareRoot{\dfrac{D+Bi}{A}}\cdot rM=qM'&\text{for even }h,\\[4pt]
{}=pM'&\text{for odd }h.
\end{array}
\]

If one derives from $h$, $h'$, $h''$ etc. the transformation of $(A,B,A')$ into $(a,b,c)$, then its elements, according as they are even or odd, determine which function of $\dfrac{D+bi}{a}$ is so re\SourceLeaf{108}lated to the given function of $M$ that
\[
\varepsilon^H\BareRoot{\left(\dfrac{D+Bi}{A}\cdot\dfrac{D+B'i}{A'}\cdot\dfrac{D+B''i}{A''}\cdots\right)}
\]
is the factor by which the latter must be multiplied.

Where $M$ is not rational, one may set $D=-1$ and construct the algorithm in the same way. Namely, if $M=g+hi$, one starts from the form $\left(\dfrac1g,\dfrac hg,\dfrac{gg+hh}{g}\right)$ (Det. $-1$), seeks its equivalent forms, etc.

\begin{center}\rule{3cm}{.4pt}\end{center}
\begin{center}[3.]\end{center}
\[
\begin{aligned}
pt&=1+2e^{-\pi t}+2e^{-4\pi t}+2e^{-9\pi t}+\text{etc.}\\
qt&=1-2e^{-\pi t}+2e^{-4\pi t}-2e^{-9\pi t}+\text{etc.}\\
rt&=2e^{-\frac14\pi t}+2e^{-\frac94\pi t}+2e^{-\frac{25}4\pi t}+\text{etc.}
\end{aligned}
\]

To solve the equation
\[
\frac{qt}{pt}=A
\]
set $AA=\dfrac nm$ and find the arithmetic-geometric mean of $m$ and $n$. Let this be $=\mu$. Then find the arithmetic-geometric mean of $m$ and $\BareRoot{(mm-nn)}$ or, equivalently, of $(m+n)$ and $(m-n)$. Let this be $=\lambda$. Then we have $t=\dfrac\mu\lambda$.

This yields only one value of $t$. All the others are contained in the formula
\[
\frac{\alpha t-2\beta i}{\delta-2\gamma ti}
\]
where $\alpha$, $\beta$, $\gamma$, $\delta$ range over all integers for which the equation
\[
\alpha\delta-4\beta\gamma=1
\]
holds.

To derive from $A$ the value
\[
B=\dfrac{q\!\left(\frac13t\right)}{p\!\left(\frac13t\right)},
\]
\SourceLeaf{109}one must solve a quartic equation,
\[
(B-A)^4=4(A-A^3)(B-B^3),
\]
or
\[
(1-AB)^4=(1-A^4)(1-B^4).
\]

The four roots correspond to $\frac13t$, $\frac13t+\frac23i$, $\frac13t+\frac43i$, $3t$.

For $A=\frac12$ find the two arithmetic-geometric means
\[
\begin{aligned}
m&=4, & n&=1\\
\mu&=2{,}2430340, &\lambda&=3{,}9364917.
\end{aligned}
\]
Hence
\[
\begin{array}{rl}
t=&0{,}56983,\\
\log t=&9{,}7557537,\\
\log(\pi\log e)=&0{,}1349342,\\\hline
&9{,}8906879,\\
\log e^{-\pi t}=&-0{,}7774777,\\
{}=&9{,}2225223,\\
pt=&1{,}33540375,\\
qt=&0{,}66770187.
\end{array}
\]
\begin{center}\rule{3cm}{.4pt}\end{center}
